\documentclass[11pt,a4paper]{article}

\usepackage[utf8]{inputenc}
\usepackage[T1]{fontenc}
\usepackage{lmodern}
\usepackage[margin=2.2cm]{geometry}
\usepackage{xcolor}
\usepackage{listings}
\usepackage{booktabs}
\usepackage{tabularx}
\usepackage{hyperref}
\usepackage{parskip}

\definecolor{kw}{RGB}{0,70,160}
\definecolor{cm}{RGB}{110,110,110}
\definecolor{st}{RGB}{160,40,40}
\definecolor{bg}{RGB}{247,247,250}

\lstdefinelanguage{Lean}{
  morekeywords={def,theorem,structure,class,instance,inductive,where,namespace,
    end,deriving,extends,by,do,let,mut,for,in,return,fun,match,with,Type,Prop},
  sensitive=true,
  morecomment=[l]{--},
  morecomment=[s]{/-}{-/},
  morestring=[b]",
}

\lstdefinestyle{base}{
  basicstyle=\ttfamily\small,
  keywordstyle=\color{kw}\bfseries,
  commentstyle=\color{cm}\itshape,
  stringstyle=\color{st},
  backgroundcolor=\color{bg},
  frame=single, rulecolor=\color{black!15},
  columns=fullflexible, keepspaces=true,
  xleftmargin=4pt, xrightmargin=4pt,
  aboveskip=6pt, belowskip=6pt,
  literate={α}{{$\alpha$}}1 {→}{{$\to$}}1 {·}{{$\cdot$}}1,
}
\lstset{style=base}

\newcommand{\lean}[1]{\lstinline[language=Lean]{#1}}

\title{Functional vs.\ Object-Oriented Programming\\[2pt]
  \large in the context of Lean~4}
\author{}
\date{}

\begin{document}
\maketitle
\vspace{-2.5em}

\section*{The short version}

Lean~4 is a \textbf{purely functional} language. It has no classes, no
objects with hidden mutable state, and no subclass polymorphism. Yet many
OOP ideas have a functional counterpart in Lean:

\begin{center}
\begin{tabularx}{\linewidth}{@{}lX@{}}
\toprule
\textbf{OOP concept} & \textbf{Lean counterpart} \\
\midrule
Class with fields        & \lean{structure} \\
Methods                  & Functions in the structure's \lean{namespace}, called with dot notation \\
Inheritance              & \lean{structure B extends A} (composition + automatic projection) \\
Interface / abstract class & Type class (\lean{class} + \lean{instance}) \\
Class hierarchy (closed) & \lean{inductive} type + pattern matching \\
Mutating a field         & Building a new value: \texttt{\{ p with x := ... \}} \\
Mutable state            & Monads (\lean{StateM}, \lean{IO.Ref}) and \lean{let mut} sugar \\
\bottomrule
\end{tabularx}
\end{center}

The key difference: in Lean every function is a \emph{mathematical}
function. Same input, same output, no side effects. That is exactly what
lets Lean \textbf{prove theorems about programs}.

%--------------------------------------------------------------------
\section{Data and behaviour}

\textbf{OOP} bundles data and methods into one class; methods act on
\texttt{this}.

\begin{lstlisting}[language=Java]
class Point {
  int x, y;
  int normSq()             { return x*x + y*y; }
  void move(int dx, int dy) { x += dx; y += dy; }   // mutates this
}
\end{lstlisting}

\textbf{Lean} separates data (\lean{structure}) from functions. Putting the
functions in the \lean{Point} namespace enables dot notation, so it
\emph{looks} like a method call, but \lean{p.normSq} is just
\lean{Point.normSq p}.

\begin{lstlisting}[language=Lean]
structure Point where
  x : Int
  y : Int

namespace Point
def normSq (p : Point) : Int := p.x * p.x + p.y * p.y
def move (p : Point) (dx dy : Int) : Point :=
  { p with x := p.x + dx, y := p.y + dy }   -- returns a NEW point
end Point

def p : Point := { x := 3, y := 4 }
#eval p.normSq         -- 25
#eval (p.move 1 1).x   -- 4
#eval p.x              -- 3   (p is unchanged)
\end{lstlisting}

%--------------------------------------------------------------------
\section{Inheritance vs.\ structure extension}

\textbf{OOP:} \texttt{class ColoredPoint extends Point}. A
\texttt{ColoredPoint} \emph{is a} \texttt{Point} and can override its
methods.

\textbf{Lean:} \lean{extends} copies the parent's fields and generates a
projection \lean{toPoint}. Dot notation searches the parent namespace, so
parent functions still work. There is \emph{no overriding} and no dynamic
dispatch: this is composition, not subtyping.

\begin{lstlisting}[language=Lean]
structure ColoredPoint extends Point where
  color : String

def cp : ColoredPoint := { x := 1, y := 2, color := "red" }
#eval cp.normSq        -- 5   (resolved as Point.normSq cp.toPoint)
\end{lstlisting}

%--------------------------------------------------------------------
\section{Polymorphism}

OOP uses one mechanism (virtual methods) for two different needs. Lean
splits them.

\subsection*{Open set of types $\to$ type classes}

Like an interface: anyone can add a new type later by writing an instance.
Dispatch is resolved at \emph{compile time} by instance search, not at run
time via a vtable.

\begin{lstlisting}[language=Lean]
class HasArea (α : Type) where
  area : α → Nat

structure Square where side : Nat
structure Rect   where (w h : Nat)

instance : HasArea Square where area s := s.side * s.side
instance : HasArea Rect   where area r := r.w * r.h

def totalArea {α : Type} [HasArea α] (xs : List α) : Nat :=
  (xs.map HasArea.area).sum

#eval totalArea [Square.mk 2, Square.mk 3]   -- 13
\end{lstlisting}

\subsection*{Closed set of cases $\to$ inductive types}

Like a sealed class hierarchy. All variants are known up front, and
pattern matching must be \emph{exhaustive}, so the compiler rejects a
forgotten case.

\begin{lstlisting}[language=Lean]
inductive Shape where
  | square (side : Nat)
  | rect   (w h : Nat)

def Shape.area : Shape → Nat
  | .square s => s * s
  | .rect w h => w * h

#eval ([.square 2, .rect 2 5] : List Shape).map Shape.area   -- [4, 10]
\end{lstlisting}

\textit{Trade-off (the ``expression problem''):} OOP makes adding a new
\emph{type} easy and a new \emph{operation} hard. Inductive types make
adding a new \emph{operation} easy (just write a function) and a new
\emph{case} hard (every match must be updated).

%--------------------------------------------------------------------
\section{State and mutation}

OOP objects carry mutable state. Lean values are immutable; state is
\emph{modelled explicitly}. The \lean{do}-notation with \lean{let mut}
looks imperative, but it compiles to pure code.

\begin{lstlisting}[language=Lean]
def sumTo (n : Nat) : Nat := Id.run do
  let mut acc := 0
  for i in [0:n+1] do
    acc := acc + i
  return acc
#eval sumTo 10                    -- 55

-- A "counter object" as an explicit state monad
def tick : StateM Nat Unit := modify (· + 1)
def threeTicks : StateM Nat Unit := do tick; tick; tick
#eval threeTicks.run 0            -- ((), 3)
\end{lstlisting}

%--------------------------------------------------------------------
\section{Why it matters: proofs}

Since functions are pure, they can be reasoned about like equations. The
same file that defines \lean{Point.move} can prove facts about it, checked
by Lean's kernel.

\begin{lstlisting}[language=Lean]
theorem move_zero (q : Point) : q.move 0 0 = q := by
  simp [Point.move]

theorem area_square (n : Nat) : (Shape.square n).area = n * n := rfl
\end{lstlisting}

In an OOP setting, \texttt{p.move(0,0)} could in principle touch any
reachable state (logging, other objects, globals), so a statement like the
one above would not even be well defined without extra assumptions.

%--------------------------------------------------------------------
\section{Summary}

\begin{center}
\begin{tabularx}{\linewidth}{@{}lXX@{}}
\toprule
 & \textbf{OOP} & \textbf{Lean (functional)} \\
\midrule
Core unit      & Object (state + behaviour) & Value + pure function \\
State          & Mutable, hidden in objects & Immutable; explicit via monads \\
Reuse          & Inheritance                & \lean{extends}, type classes, composition \\
Dispatch       & Dynamic (run time)         & Static (instances) or pattern matching \\
Extensibility  & Easy new types             & Easy new operations (inductives) \\
Correctness    & Tests                      & Tests \emph{and} machine-checked proofs \\
\bottomrule
\end{tabularx}
\end{center}

\vfill
\small\textit{All Lean snippets typecheck with Lean 4.35.0 (core only, no Mathlib).}

\end{document}
